\section{Принципы нейросетевого сжатия изображений на базе вариационных автокодировщиков}

\subsection{Общая постановка задачи сжатия изображений}
Сжатие изображений с потерями представляет собой фундаментальную задачу теории информации и компьютерного зрения, направленную на максимально компактное представление визуальных данных при допустимом уровне искажений. Традиционные методы (JPEG, WebP) основаны на ручном проектировании линейных преобразований (например, дискретное косинусное преобразование — DCT) и статистических моделей энтропии кодирования. Однако эти подходы достигают предела своей эффективности, так как не адаптируются к специфике конкретных изображений и не используют семантические особенности данных.

Нейросетевое сжатие (learned image compression) преодолевает эти ограничения путем совместной оптимизации всех этапов сжатия — анализа, квантования и энтропийного кодирования — с помощью глубоких нейронных сетей. Ключевыми преимуществами такого подхода являются:
\begin{itemize}
\item Адаптивность: модели обучаются на больших наборах данных и учатся выделять наиболее значимые для восприятия и реконструкции признаки.
\item Нелинейные преобразования: нейросети способны реализовывать сложные нелинейные преобразования, более эффективные, чем фиксированные линейные.
\item Сквозная оптимизация: вся система настраивается для прямой минимизации целевой функции, балансирующей между размером сжатого представления (битрейтом) и качеством реконструкции.
\item Поддержка новых форматов: возможность включения в конвейер сжатия перцептивных или задачно-ориентированных функций потерь.
\end{itemize}
Таким образом, нейросетевое сжатие открывает новый этап в развитии технологий компрессии, потенциально превосходящий традиционные стандарты на границах компромисса «битрейт-качество»~\cite{balle2016, minnen2018}.

\subsection{Основы вариационных автокодировщиков (VAE)}
Вариационный автокодировщик (Variational Autoencoder, VAE) — это генеративная вероятностная модель, которая позволяет обучать сложные распределения данных в рамках байесовского подхода. В контексте сжатия VAE рассматривается как система, состоящая из двух основных компонентов:
\begin{itemize}
\item \textbf{Энкодер (анализ):} Преобразует входное изображение $\mathbf{x}$ в параметры ($\mathbf{\mu}$, $\mathbf{\sigma}$) компактного латентного представления $\mathbf{z}$, имеющего вероятностную природу.
\item \textbf{Декодер (синтез):} Восстанавливает изображение $\mathbf{\hat{x}}$ из латентной переменной $\mathbf{z}$.
\end{itemize}
Принципиальное отличие VAE от классического автокодировщика заключается в трактовке латентного пространства. В VAE накладывается априорное распределение (обычно нормальное) $p(\mathbf{z})$, а энкодер предсказывает параметры апостериорного распределения $q(\mathbf{z}|\mathbf{x})$. Обучение модели происходит путем максимизации нижней оценки (Evidence Lower BOund, ELBO) правдоподобия данных:

\begin{equation}\label{eq:elbo}
    \log p(\mathbf{x}) \geq \mathrm{E}_{q(\mathbf{z}|\mathbf{x})}[\log p(\mathbf{x}|\mathbf{z})] - D_{\mathrm{KL}}(q(\mathbf{z}|\mathbf{x}) \| p(\mathbf{z})).
\end{equation}

Данная формула напрямую связывает VAE с задачей сжатия: первый член (правдоподобие) отвечает за качество реконструкции (искажение, \textit{distortion}), а второй член (KL-дивергенция) --- за сложность латентного представления (скорость, \textit{rate}). Минимизация KL-дивергенции приближает апостериорное распределение к простому априорному, что облегчает эффективное энтропийное кодирование~\cite{kingma2013}.

\subsection{Принцип работы VAE для сжатия изображений}
Архитектура VAE для сжатия изображений является прямой реализацией описанных выше теоретических принципов в практический конвейер кодирования и декодирования. Его работа включает следующие ключевые этапы:

\begin{itemize}
\item \textbf{Анализ (Кодирование):} Энкодер $f_e$ преобразует изображение $\mathbf{x}$ в латентное представление $\mathbf{y} = f_e(\mathbf{x}; \phi)$, которое затем подвергается квантованию $\mathbf{\hat{y}} = Q(\mathbf{y})$ для получения дискретного представления, пригодного для энтропийного кодирования.
\item \textbf{Энтропийное кодирование:} Квантованные латенты $\mathbf{\hat{y}}$ сжимаются арифметическим кодером, использующим заранее обученную вероятностную модель (это и есть априорное распределение $p_{\mathbf{\hat{y}}}$). Размер битстрима определяется энтропией этого распределения.
\item \textbf{Синтез (Декодирование):} Декодер $f_d$ восстанавливает изображение $\mathbf{\hat{x}} = f_d(\mathbf{\hat{y}}; \theta)$ из деквантованных латентов.
\end{itemize}

Обучение такой системы представляет собой задачу оптимизации компромисса между скоростью (rate) и искажениями (distortion), что формализуется минимизацией функции потерь вида:

\begin{equation}\label{eq:rd_loss}
    \mathcal{L} = \lambda \cdot D(\mathbf{x}, \mathbf{\hat{x}}) + R(\mathbf{\hat{y}}),
\end{equation}
где $D$ — функция искажения (например, MSE или MS-SSIM), $R$ — оценка битрейта (отрицательный логарифм вероятности латентов), а $\lambda$ — параметр, контролирующий баланс между ними. Для обеспечения дифференцируемости процесса квантования во время обучения используется метод добавления униформного шума (uniform noise) как его дифференцируемая аппроксимация~\cite{balle2016}.

Таким образом, применение принципов VAE позволяет создать полностью дифференцируемую систему сжатия, которая напрямую оптимизирует фундаментальный компромисс rate-distortion и способна превзойти традиционные методы за счет адаптивного нелинейного анализа данных.